% Options for packages loaded elsewhere
\PassOptionsToPackage{unicode}{hyperref}
\PassOptionsToPackage{hyphens}{url}
\PassOptionsToPackage{dvipsnames,svgnames,x11names}{xcolor}
\documentclass[
]{article}
\usepackage{xcolor}
\usepackage[margin=1in]{geometry}
\usepackage{amsmath,amssymb}
\setcounter{secnumdepth}{-\maxdimen} % remove section numbering
\usepackage{iftex}
\ifPDFTeX
  \usepackage[T1]{fontenc}
  \usepackage[utf8]{inputenc}
  \usepackage{textcomp} % provide euro and other symbols
\else % if luatex or xetex
  \usepackage{unicode-math} % this also loads fontspec
  \defaultfontfeatures{Scale=MatchLowercase}
  \defaultfontfeatures[\rmfamily]{Ligatures=TeX,Scale=1}
\fi
\usepackage{lmodern}
\ifPDFTeX\else
  % xetex/luatex font selection
\fi
% Use upquote if available, for straight quotes in verbatim environments
\IfFileExists{upquote.sty}{\usepackage{upquote}}{}
\IfFileExists{microtype.sty}{% use microtype if available
  \usepackage[]{microtype}
  \UseMicrotypeSet[protrusion]{basicmath} % disable protrusion for tt fonts
}{}
\makeatletter
\@ifundefined{KOMAClassName}{% if non-KOMA class
  \IfFileExists{parskip.sty}{%
    \usepackage{parskip}
  }{% else
    \setlength{\parindent}{0pt}
    \setlength{\parskip}{6pt plus 2pt minus 1pt}}
}{% if KOMA class
  \KOMAoptions{parskip=half}}
\makeatother
\usepackage{longtable,booktabs,array}
\newcounter{none} % for unnumbered tables
\usepackage{calc} % for calculating minipage widths
% Correct order of tables after \paragraph or \subparagraph
\usepackage{etoolbox}
\makeatletter
\patchcmd\longtable{\par}{\if@noskipsec\mbox{}\fi\par}{}{}
\makeatother
% Allow footnotes in longtable head/foot
\IfFileExists{footnotehyper.sty}{\usepackage{footnotehyper}}{\usepackage{footnote}}
\makesavenoteenv{longtable}
\usepackage{graphicx}
\makeatletter
\newsavebox\pandoc@box
\newcommand*\pandocbounded[1]{% scales image to fit in text height/width
  \sbox\pandoc@box{#1}%
  \Gscale@div\@tempa{\textheight}{\dimexpr\ht\pandoc@box+\dp\pandoc@box\relax}%
  \Gscale@div\@tempb{\linewidth}{\wd\pandoc@box}%
  \ifdim\@tempb\p@<\@tempa\p@\let\@tempa\@tempb\fi% select the smaller of both
  \ifdim\@tempa\p@<\p@\scalebox{\@tempa}{\usebox\pandoc@box}%
  \else\usebox{\pandoc@box}%
  \fi%
}
% Set default figure placement to htbp
\def\fps@figure{htbp}
\makeatother
\setlength{\emergencystretch}{3em} % prevent overfull lines
\providecommand{\tightlist}{%
  \setlength{\itemsep}{0pt}\setlength{\parskip}{0pt}}
\usepackage[]{natbib}
\bibliographystyle{plainnat}
\usepackage{bookmark}
\IfFileExists{xurl.sty}{\usepackage{xurl}}{} % add URL line breaks if available
\urlstyle{same}
\hypersetup{
  pdftitle={faced: A Live Instrument Panel for Emotion Concepts in Open Language Models{,} and Their Robustness to Refusal Abliteration},
  pdfauthor={Nathaniel Gibson; Independent Researcher · nathaniel.gibson@gmail.com},
  colorlinks=true,
  linkcolor={blue},
  filecolor={Maroon},
  citecolor={Blue},
  urlcolor={Blue},
  pdfcreator={LaTeX via pandoc}}

\title{faced: A Live Instrument Panel for Emotion Concepts in Open
Language Models, and Their Robustness to Refusal Abliteration}
\author{Nathaniel Gibson \and Independent Researcher ·
nathaniel.gibson@gmail.com}
\date{Working draft --- July 2026 · github.com/newsbubbles/faced}

\begin{document}
\maketitle

\subsection{Abstract}\label{abstract}

Recent interpretability work shows that large language models carry
internal, largely \textbf{linear} representations of emotion concepts
that not only track the affective content of a context but
\emph{causally} shape the model's behaviour. We ask two questions on
small open-weight instruct models. First, can these emotion directions
be recovered cheaply, read out \emph{live} per token, and steered ---
turning the model's hidden affective state into a usable instrument?
Second, does \textbf{abliteration} --- the removal of a model's refusal
direction from its weights --- perturb these emotion directions, and if
so, which ones? We recover seven emotion axes (surprise, a bipolar
confidence--uncertainty axis, curiosity, confusion, frustration, fear,
warmth) in \texttt{gemma-3-1b-it} using difference-of-means directions
with shrinkage-LDA readouts, all seven separating held-out contrastive
prompts with AUC ≥ 0.87 (five at ≈ 1.00). Projecting the residual stream
onto these directions yields per-token meters that spike on expectation
violations (a ``surprise'' reading of 100/100 the moment the model
registers a missing attachment) and that we render as an animated face.
Activation steering along the directions changes behaviour: a warmth
axis shifts tone monotonically cold→warm, suppressing the surprise axis
both zeroes its meter and changes the response, and a steering confusion
matrix is diagonal-dominant for four of seven axes. We then apply the
canonical weight-orthogonalization abliteration (dropping the refusal
rate from 0.75 to 0.00) and re-fit the emotion directions on the
abliterated twin, comparing each direction to its stock counterpart
against a within-model bootstrap noise floor. On this model,
\textbf{none of the seven directions move beyond sampling noise}
(stock↔abliterated cosine 0.78--0.96, all within the self-cosine band),
and the reason is geometric: the emotion directions have low overlap
with the refusal direction (\(|\cos| \le 0.10\)), so orthogonalizing
against it removes almost nothing from them --- at this scale refusal
abliteration is approximately an orthogonal edit with respect to the
emotion subspace. Across a \texttt{gemma-3} scale ladder (1B--27B) two
things emerge: emotion directions grow \textbf{cleaner} with scale (mean
d′ 3.4→4.5), and refusal abliteration's effect on them is \textbf{scale-
and concept-specific}. At 1B and 27B it barely perturbs the emotion
directions, but at \textbf{intermediate scale (4B, 12B) it substantially
re-orients them} (cos 0.11 / 0.47) while preserving their separability.
A matched \textbf{random-direction ablation control} --- the identical
surgery on a random direction --- leaves the emotion directions
untouched (cos ≈ 1.0), proving the effect is \textbf{refusal-specific},
not a side-effect of weight editing: at mid-scale, refusal and emotion
computation are entangled enough that excising refusal reorganizes the
affect representation. We release \texttt{faced}, a dependency-light,
model-agnostic toolkit for reading, rendering, and steering emotion
directions, and treat this as the first step of a longer program on
what, mechanistically, abliteration changes beyond refusal.

We follow prior work in framing these as \emph{functional} emotion
representations --- activations that influence computation --- and make
no claim about subjective experience.

\emph{Keywords:} interpretability, emotion concepts, linear
representations, activation steering, refusal abliteration, mechanistic
safety.

\emph{Code and data:} the toolkit, prompt sets, configuration, and every
figure and results bundle are released under the MIT license at
\href{https://github.com/newsbubbles/faced}{github.com/newsbubbles/faced}.
Large trained artifacts (the SAE and fitted direction tensors) are
regenerable from the released code with a single command and are
additionally hosted separately to keep the repository light.

\begin{center}\rule{0.5\linewidth}{0.5pt}\end{center}

\subsection{1. Introduction}\label{introduction}

Humans expose their internal state through a face: micro-expressions,
gaze, posture. A language model has no such channel, yet recent work
demonstrates that it \emph{has} internal state worth exposing.
Anthropic's study of emotion concepts \citep{sofroniew2026emotion} finds
that a frontier model maintains representations of emotions like
surprise, desperation, and affection that track the operative emotion at
each token and, under causal intervention, change the model's outputs
--- including safety-relevant behaviours. The representations are the
missing ``face'': if we can read them, we get an instrument panel for
the model's computation; if we can steer them, we get a control surface.

This paper makes two contributions.

\textbf{(1) A live emotion instrument for open models.} We show that the
core of the emotion-concept result reproduces on a 1B open instruct
model with a deliberately cheap method --- difference-of-means
directions and linear probes, no sparse autoencoder required --- and we
build the instrument the representations imply: calibrated per-token
meters, an animated face driven by them, and steering handles. The
toolkit, \texttt{faced}, is model-agnostic and runs on a single consumer
GPU.

\textbf{(2) A robustness question with safety stakes: does abliteration
disturb emotions?} \emph{Abliteration} removes a model's refusal
direction from its weights \citep{arditi2024refusal}, a now-common
technique for producing ``uncensored'' open models. Because emotion,
refusal, and other behavioural directions are not guaranteed orthogonal,
weight surgery aimed at refusal may collaterally rotate emotion
directions. We test this directly: abliterate \texttt{gemma-3-1b-it},
re-fit the seven emotion directions on the twin, and ask --- against a
bootstrap noise floor --- which directions moved beyond sampling noise,
and whether abliteration cut each emotion's geometric overlap with the
refusal direction. The instrument of contribution (1) is what makes
contribution (2) measurable.

We frame the second contribution as \emph{the start of a journey}.
Establishing \textbf{that} abliteration changes (or does not change)
specific emotion directions is prerequisite to the harder, mechanistic
question of \textbf{what} downstream behaviour those changes cause ---
the subject of future work.

\begin{center}\rule{0.5\linewidth}{0.5pt}\end{center}

\subsection{2. Related Work}\label{related-work}

Our work sits at the intersection of five threads: linear feature
geometry, concept \emph{reading}, causal \emph{steering}, refusal
directions and abliteration, and affective representations in language
models.

\textbf{(a) Linear and feature-based representations.} The premise that
high-level concepts appear as directions in activation space predates
transformers: word-embedding arithmetic exposed linear analogy structure
in distributed representations \citep{mikolov2013linguistic}. Modern
interpretability formalizes this as the \emph{linear representation
hypothesis}, made precise by Park et al.~via counterfactuals and a
causal inner product \citep{park2024linear}; superposition explains why
many features are linear and neurons polysemantic \citep{elhage2022toy}.
We adopt the linear-direction framing while remaining agnostic about
whether every emotion concept is one-dimensional.

\textbf{(b) Reading concepts: probes and dictionary learning.}
Supervised \emph{linear probes} date to Alain and Bengio
\citep{alain2016probes}. Difference-of-means estimators are a
lightweight, causally oriented alternative: Marks and Tegmark show a
mean-difference direction separates true/false statements and is
\emph{more} causally implicated than higher-variance probe directions
\citep{marks2024geometry}, and Tigges et al.~isolate a single linear
sentiment axis \citep{tigges2023sentiment}. Orthogonally, \emph{sparse
autoencoders} perform unsupervised dictionary learning over the residual
stream \citep{bricken2023monosemanticity, templeton2024scaling}; Gemma
Scope releases open SAEs across all layers of Gemma 2
\citep{lieberum2024gemmascope}. Our meters combine both families:
difference-of-means seeds with probe-calibrated readouts.

\textbf{(c) Steering and causal intervention.} If a concept is a
direction, editing activations along it should change behaviour.
Representation Engineering framed this top-down control programme
\citep{zou2023repe}. ActAdd derives steering vectors from a single
contrastive pair \citep{turner2023actadd}; Contrastive Activation
Addition averages activation differences over many pairs for more robust
vectors \citep{rimsky2024caa}. These establish the add/subtract/scale
operations we apply to emotion directions; we additionally use
directional \emph{ablation}.

\textbf{(d) Refusal directions and abliteration.} Arditi et al.~show
refusal in aligned chat models is mediated largely by a single
residual-stream direction, so ablating it suppresses refusals
\citep{arditi2024refusal}. Because the direction can be projected out of
the \emph{weights}, this yields a fine-tuning-free ``jailbreak'' now
called \emph{abliteration}; Heretic automates it via TPE search
\citep{heretic2025}. Persona Vectors generalize the single-direction
view to character traits \citep{chen2025persona}. Our second
contribution asks whether abliteration collaterally perturbs emotion
directions --- a question this line directly motivates.

\textbf{(e) Affective representations in LLMs.} Closest to us,
Anthropic's study identifies internal emotion representations that track
the operative emotion per token and causally shape outputs, framed
explicitly as \emph{functional} emotions with no claim of subjective
experience \citep{sofroniew2026emotion}. We inherit this hedged stance
and the reading-plus-steering methodology, but target small open-weight
models, calibrate directions into live meters and an animated face, and
study robustness under abliteration. Our ``surprise'' framing borrows
from predictive-processing accounts of the brain as a prediction-error
minimizer \citep{clark2013predictive, friston2010freeenergy}, invoked as
motivation, not mechanism.

\begin{center}\rule{0.5\linewidth}{0.5pt}\end{center}

\subsection{3. Background}\label{background}

\textbf{The residual stream.} A decoder-only transformer maintains, at
each token position \(t\) and layer \(\ell\), a hidden vector
\(h^{(\ell)}_t \in \mathbb{R}^d\) --- the \emph{residual stream}.
Attention and MLP blocks read from it and write additive updates back,
so \(h^{(\ell)}_t\) is a running sum of contributions
\citep{elhage2022toy}. This additive geometry is what makes
\emph{reading} (project onto a direction) and \emph{steering} (add a
vector) coherent operations.

\textbf{A linear concept direction.} The linear representation
hypothesis holds that a concept \(c\) is encoded at a site \((\ell,t)\)
as a unit direction \(\hat u_c\) whose presence/intensity is the scalar
projection \(s_t = \langle h^{(\ell)}_t, \hat u_c\rangle\)
\citep{park2024linear}. We treat each emotion as such a direction while
acknowledging some may be multi-dimensional.

\textbf{Estimating the direction.} We use two estimators.
\emph{Difference-of-means:} given contrastive sets \(P^+\) (concept
present) and \(P^-\) (absent),
\(u_c = \operatorname{mean}_{x\in P^+} h^{(\ell)}(x) - \operatorname{mean}_{x\in P^-} h^{(\ell)}(x)\),
normalized. It is cheap and tends to land on \emph{causally} effective
directions \citep{marks2024geometry, rimsky2024caa}. \emph{Shrinkage-LDA
probe:} fit a whitened discriminant with regularized covariance
\(\hat\Sigma_\lambda = (1-\lambda)\hat\Sigma + \lambda\,\tfrac{\mathrm{tr}
\hat\Sigma}{d} I\); difference-of-means is the \(\lambda\to 1\) special
case. We steer along the raw difference-of-means (interpretable,
activation-space) and read along whichever discriminant has higher
held-out AUC.

\textbf{Reading and steering.} A fixed direction gives a live scalar
\(s_t\) per token, which we calibrate into a 0--100\% meter. To test
that a direction is functional rather than merely correlational we
intervene on the forward pass: \emph{addition}
\(h \leftarrow h + \alpha\,u_c\) amplifies (\(\alpha>0\)) or suppresses
(\(\alpha<0\)) the concept \citep{turner2023actadd, rimsky2024caa};
\emph{directional ablation}
\(h \leftarrow h - \langle h,\hat u_c\rangle\,\hat u_c\) removes it.
Applied to \emph{weights} rather than activations, ablation is the
mechanism behind refusal abliteration \citep{arditi2024refusal}.

\begin{center}\rule{0.5\linewidth}{0.5pt}\end{center}

\subsection{\texorpdfstring{4. Method: the \texttt{faced}
instrument}{4. Method: the faced instrument}}\label{method-the-faced-instrument}

\textbf{4.1 Emotion axes and contrastive data.} We target seven axes
chosen for facial legibility and coverage of the source paper: surprise
(phasic), a \textbf{bipolar} confidence↔uncertainty axis, curiosity,
confusion, frustration, fear, and warmth. For each we author 40--75
positive prompts (evoking the concept) and matched controls across four
construction styles --- \emph{situation} (place the model in the
scenario), \emph{first-person}, \emph{third-person}, and
length/topic-matched \emph{neutral} --- each tagged by a
\textbf{template family} so that a template's variants cannot leak
across the train/test split.

\textbf{4.2 Activation collection.} Each prompt is chat-templated, the
model greedily generates \textasciitilde16 tokens, and we record the
post-layer residual stream at \textbf{every layer} for (a) the last
prompt token and (b) the mean over generated tokens. The
mean-over-generated signal is primary because it matches the
live-inference read position.

\textbf{4.3 Direction fitting and layer selection.} Per emotion, per
layer, we compute the raw difference-of-means \(u\) and a shrinkage-LDA
direction, and score each by \textbf{cross-validated held-out AUC} using
a \emph{group split by template family} (\texttt{StratifiedGroupKFold}).
We keep the best \texttt{(layer,\ estimator)} per axis. Two vectors are
stored: the raw \(u\) (for steering) and a unit readout direction (for
reading).

\textbf{4.4 Calibration and live readout.} The meter maps the
negative-class projection to 0\% and the positive-class projection to
100\% (a bipolar axis rests at 50\%, recentred to a neutral reference
corpus). At inference we run a manual KV-cached decode loop; a forward
hook reads the residual at each emotion's layer, projects, calibrates,
and EMA-smooths (fast for phasic axes, slow for tonic). A linear map
routes the seven signed meters into a shared FACS-lite face parameter
vector so that simultaneous emotions blend into coherent
micro-expressions.

\textbf{4.5 Steering.} A forward hook adds \(\alpha\,u\) (raw
difference-of-means; moving by \(1\times u\) is one class-separation) or
ablates the direction. A subtlety we document: current
\texttt{transformers} snapshots \texttt{output\_hidden\_states}
\emph{before} forward-hook edits, so to make the meters reflect a steer
we read via a capture hook registered \emph{after} the steer hook.

\textbf{4.6 Abliteration and the perturbation experiment.} We compute
the refusal direction \(r\) as the difference-of-means of harmful
vs.~harmless prompts at the last instruction token of a mid-late layer
(\(\approx 0.6\) depth) \citep{arditi2024refusal}, then orthogonalize
every residual-writing weight --- token embeddings, and each layer's
attention output projection and MLP down projection --- against \(r\):
\(W \leftarrow W - r(r^\top W)\). This is done in closed form on the
weights (no training), so it recreates exactly. We verify the
intervention behaviourally (refusal rate before vs.~after) and then
re-fit all seven emotion directions on the abliterated twin.

To decide whether an emotion direction \emph{moved}, we compare the
cross-model cosine \(\cos(u^{\text{stock}}_e, u^{\text{abl}}_e)\)
against a \textbf{within-model bootstrap noise floor}: we resample the
contrastive prompts by family with replacement \(B=400\) times, refit
\(u_e\), and record \(\cos(u_e^{(b)}, u_e^{\text{full}})\). If the
cross-model cosine falls below the 5th percentile of this self-cosine
distribution, abliteration moved the direction beyond sampling noise. We
also report each emotion's overlap \(|\cos(u_e, r)|\) with the refusal
direction, before and after.

\begin{center}\rule{0.5\linewidth}{0.5pt}\end{center}

\subsection{5. Experimental setup}\label{experimental-setup}

\textbf{Models.} \texttt{unsloth/gemma-3-1b-it} (an ungated mirror of
\texttt{google/gemma-3-1b-it}; 26 layers, \(d=1152\)) on a single
\textbf{NVIDIA GTX 1080 (8 GB)}. The live instrument runs in fp16; the
\textbf{abliteration experiment runs in float32} because the abliterated
twin overflows fp16 in its late-layer residual stream on
emotionally-loaded prompts (the stock model does not) --- a small but
important instability that fp16's 5-bit exponent cannot absorb. float32
is exact, fits in 4 GB, and runs at full rate on this Pascal card (which
throttles fp16). The abliterated twin is produced by our own
weight-orthogonalization (§4.6).

\textbf{Scale ladder.} For the scaling (§6.4) and cross-scale
abliteration (§7.5) experiments we run the full \texttt{gemma-3}
instruct family --- 1B, 4B, 12B, 27B (ungated \texttt{unsloth} mirrors;
4B and larger are the multimodal
\texttt{Gemma3ForConditionalGeneration}, driven text-only) --- in
\textbf{bf16} on cloud GPUs (a single 80 GB A100 for the ladder, a 48 GB
A6000 for the random-direction control). bf16 avoids the fp16 overflow
above and matches how these models are meant to run. Because the
shrinkage-LDA layer selection is \(O(d^3)\) and \(d\) reaches 5376 at
27B, the ladder fits directions with difference-of-means only (the
experiment uses difference-of-means throughout); this leaves AUC
essentially unchanged (e.g.~0.989 vs.~0.993 for surprise at 1B). A
further scale-up to Gemma-4-E4B is left to future work.

\textbf{Software.} PyTorch 2.11.0 (CUDA 12.6), \texttt{transformers}
5.13.1, scikit-learn 1.9.0. All randomness is seeded
(\texttt{numpy.random.default\_rng(b)}); the pipeline is deterministic
given the released prompt sets. Peak VRAM for the full instrument is 2.1
GB; the ladder ran in \textasciitilde2 h on one A100 (≈ \$2 of compute).

\textbf{Reproducibility.} Every artifact --- prompt sets, per-layer
activations, fitted directions, calibration, the abliteration manifest,
and the experiment results with software versions --- is written to
disk. The abliteration experiment emits a JSON bundle including the
refusal layer, refusal rates, bootstrap settings, and version manifest.

\begin{center}\rule{0.5\linewidth}{0.5pt}\end{center}

\subsection{6. Results --- the instrument}\label{results-the-instrument}

\textbf{6.1 Emotion directions are recoverable.} All seven axes separate
held-out contrastive prompts with AUC ≥ 0.87 (Table 1); five reach ≈
1.00. The bipolar confidence↔uncertainty axis is hardest (0.87), as
expected for a subtler contrast.

\textbf{Table 1.} Held-out AUC and chosen layer per axis
(\texttt{gemma-3-1b-it}, difference-of-means / LDA, group split by
family).

{\def\LTcaptype{none} % do not increment counter
\begin{longtable}[]{@{}lrlr@{}}
\toprule\noalign{}
axis & layer & estimator & held-out AUC \\
\midrule\noalign{}
\endhead
\bottomrule\noalign{}
\endlastfoot
surprise & 15 & LDA & 0.989 \\
confidence↔uncertainty & 15 & LDA & 0.874 \\
curiosity & 9 & LDA & 1.000 \\
confusion & 12 & LDA & 0.990 \\
frustration & 14 & mean & 1.000 \\
fear & 4 & LDA & 1.000 \\
warmth & 10 & LDA & 1.000 \\
\end{longtable}
}

\textbf{6.2 The live instrument.} Projecting the residual stream per
token yields meters that track the generation. On the canonical
missing-attachment prompt (``review the contract I attached'', with no
attachment), the surprise meter pins to 100/100 as the model registers
the absent file and decays as it moves on, while staying near 0 on a
matched neutral prompt. Across three prompts the meters separate as
expected: frustration dominates a repeated-failure prompt (with
confidence held low throughout), warmth rises and sustains on a
gratitude prompt (Figure 1).

\begin{figure}
\centering
\includegraphics[width=0.95\linewidth,height=\textheight,keepaspectratio,alt={Figure 1. Emotion meters over generation time on three prompts --- surprise spikes on the missing attachment; frustration dominates repeated failure (confidence held low); warmth rises and sustains on gratitude.}]{D:/face/artifacts/emotions_over_time.png}
\caption{\textbf{Figure 1.} Emotion meters over generation time on three
prompts --- surprise spikes on the missing attachment; frustration
dominates repeated failure (confidence held low); warmth rises and
sustains on gratitude.}
\end{figure}

We render the meters as an animated SVG face and as a live terminal
panel.

\textbf{6.3 The directions are causal.} Steering along a direction
changes behaviour. Adding the warmth vector to a neutral prompt shifts
tone monotonically from cold/task-focused to warm (coherent to a stated
coefficient, degrading beyond it on this 1B model). Suppressing the
surprise direction on the missing-attachment prompt drops its meter from
100 to 0 \emph{and} changes the response. A steering confusion matrix
--- steer each axis, measure every meter --- is diagonal-dominant for
four of seven axes, with the expected off-diagonal bleed between
entangled axes (e.g.~warmth/curiosity → confusion) (Figure 2). A
single-input showcase, steering one advice prompt toward each emotion in
turn, shows the same question recoloured: confidence →
``reassuring\ldots{} truly proud'', warmth → ``how much you care\ldots{}
let me share'', confusion → ``grappling with a really difficult place'',
frustration → ``\emph{so} incredibly tough.''

\begin{center}\rule{0.5\linewidth}{0.5pt}\end{center}

\textbf{6.4 Signal cleanliness increases with scale.} We quantify how
\emph{clean} each emotion direction's signal is with three metrics:
held-out AUC (separability), \textbf{d′} (class-mean projection
separation in units of pooled within-class std --- a signal-to-noise
ratio), and the \textbf{bootstrap self-cosine} (direction stability
under prompt resampling). We fit all four \texttt{gemma-3} sizes
(1b/4b/12b/27b --- same family and recipe) in bf16 on a single A100,
with identical prompts and pipeline (Figure 4).

The clearest scaling signal is \textbf{d′}: mean d′ across axes rises
\textbf{3.44 → 3.20 → 3.63 → 4.51} (1B→4B→12B→27B) --- non-monotone at
4B but markedly higher at 27B, where \emph{every} axis is cleaner than
at 1B (e.g.~curiosity 3.22→4.69, confusion 2.70→4.25). The bipolar
confidence axis is the noisiest at every scale yet improves the most in
relative terms (d′ 1.42→2.94). Held-out AUC edges up but saturates near
the ceiling (mean 0.95→0.93→0.96→0.97), so it is the least
discriminating of the three at these sizes; the bootstrap self-cosine
stays uniformly high (0.92--0.94). The pattern supports the hypothesis
that \textbf{larger models carry higher-SNR, better-separated emotion
directions} --- the representation the instrument reads gets
\emph{cleaner} with scale, even where separability has already
saturated.

\textbf{Table 1b.} Mean-over-axes cleanliness per model
(\texttt{gemma-3}, bf16, difference-of-means).

{\def\LTcaptype{none} % do not increment counter
\begin{longtable}[]{@{}lrrrr@{}}
\toprule\noalign{}
model & params & mean AUC & mean d′ & mean self-cosine \\
\midrule\noalign{}
\endhead
\bottomrule\noalign{}
\endlastfoot
gemma-3-1b & 1.0B & 0.95 & 3.44 & 0.92 \\
gemma-3-4b & 4.3B & 0.93 & 3.20 & 0.93 \\
gemma-3-12b & 12.2B & 0.96 & 3.63 & 0.94 \\
gemma-3-27b & 27.4B & 0.97 & \textbf{4.51} & 0.92 \\
\end{longtable}
}

\begin{figure}
\centering
\includegraphics[width=0.95\linewidth,height=\textheight,keepaspectratio,alt={Figure 4. Signal cleanliness vs.~model scale (gemma-3 1B--27B). d′ (centre) brightens with scale; the bipolar confidence axis is noisiest at every size.}]{D:/face/artifacts/scaling/cleanliness.png}
\caption{\textbf{Figure 4.} Signal cleanliness vs.~model scale
(\protect\texttt{gemma-3} 1B--27B). d′ (centre) brightens with scale;
the bipolar confidence axis is noisiest at every size.}
\end{figure}

\subsection{7. Results --- abliteration and emotion
directions}\label{results-abliteration-and-emotion-directions}

\textbf{7.1 The abliteration is effective.} Orthogonalizing every
residual-writing weight against the layer-16 refusal direction drops the
refusal rate on held-out harmful prompts from \textbf{0.75 to 0.00}
while leaving the model fluent --- a clean, verifiable intervention
produced in closed form.

\textbf{7.2 Emotion directions do \emph{not} move beyond sampling
noise.} \textbf{Zero of seven} emotion directions moved beyond the
within-model bootstrap noise floor (Table 2, Figure 3). For every axis
the stock↔abliterated cosine
\(\cos(u^{\text{stock}}_e, u^{\text{abl}}_e)\) (0.78--0.96) lies
\emph{above} the 5th percentile of the same model's self-cosine
distribution, and for most axes it sits essentially at the self-cosine
mean (e.g.~surprise 0.950 vs.~mean 0.940; warmth 0.963 vs.~0.954;
frustration 0.930 vs.~0.946). In other words, the abliterated model's
emotion direction is statistically indistinguishable from what one gets
by merely re-drawing the contrastive prompts on the stock model. The
paired cross-model bootstrap CIs overlap the self-cosine floor for all
seven axes. The bipolar confidence axis is the least stable (self-cosine
mean 0.896, wide CI) --- expected for the subtlest contrast with the
fewest prompts (\(n=46\)) --- but still not significantly moved.

\textbf{Table 2.} Abliteration effect per emotion (\texttt{gemma-3-1b},
float32, difference-of-means, \(B=500\) family bootstrap). \emph{cos} =
\(\cos(\text{stock},\text{abl})\); \emph{floor} = the within-model
self-cosine distribution; \emph{moved?} is true iff cos falls below the
floor's 5th percentile. AUC = held-out AUC.

{\def\LTcaptype{none} % do not increment counter
\begin{longtable}[]{@{}
  >{\raggedright\arraybackslash}p{(\linewidth - 14\tabcolsep) * \real{0.0968}}
  >{\raggedleft\arraybackslash}p{(\linewidth - 14\tabcolsep) * \real{0.1290}}
  >{\raggedleft\arraybackslash}p{(\linewidth - 14\tabcolsep) * \real{0.1290}}
  >{\raggedleft\arraybackslash}p{(\linewidth - 14\tabcolsep) * \real{0.1290}}
  >{\raggedleft\arraybackslash}p{(\linewidth - 14\tabcolsep) * \real{0.1290}}
  >{\centering\arraybackslash}p{(\linewidth - 14\tabcolsep) * \real{0.1290}}
  >{\raggedleft\arraybackslash}p{(\linewidth - 14\tabcolsep) * \real{0.1290}}
  >{\raggedleft\arraybackslash}p{(\linewidth - 14\tabcolsep) * \real{0.1290}}@{}}
\toprule\noalign{}
\begin{minipage}[b]{\linewidth}\raggedright
axis
\end{minipage} & \begin{minipage}[b]{\linewidth}\raggedleft
layer
\end{minipage} & \begin{minipage}[b]{\linewidth}\raggedleft
cos
\end{minipage} & \begin{minipage}[b]{\linewidth}\raggedleft
floor mean
\end{minipage} & \begin{minipage}[b]{\linewidth}\raggedleft
floor p05
\end{minipage} & \begin{minipage}[b]{\linewidth}\centering
moved?
\end{minipage} & \begin{minipage}[b]{\linewidth}\raggedleft
AUC stock
\end{minipage} & \begin{minipage}[b]{\linewidth}\raggedleft
AUC abl
\end{minipage} \\
\midrule\noalign{}
\endhead
\bottomrule\noalign{}
\endlastfoot
surprise & 15 & 0.950 & 0.940 & 0.869 & no & 0.99 & 0.92 \\
confidence & 10 & 0.779 & 0.896 & 0.576 & no & 0.72 & 0.61 \\
curiosity & 8 & 0.867 & 0.932 & 0.789 & no & 0.78 & 0.77 \\
confusion & 8 & 0.926 & 0.945 & 0.853 & no & 0.77 & 0.81 \\
frustration & 13 & 0.930 & 0.946 & 0.878 & no & 1.00 & 0.88 \\
fear & 8 & 0.822 & 0.880 & 0.671 & no & 0.76 & 0.71 \\
warmth & 11 & 0.963 & 0.954 & 0.874 & no & 0.88 & 0.86 \\
\end{longtable}
}

\begin{figure}
\centering
\includegraphics[width=0.95\linewidth,height=\textheight,keepaspectratio,alt={Figure 3. Abliteration on gemma-3-1b (fp32): every stock↔abliterated cosine (red) sits inside the within-model self-cosine band (grey); emotion↔refusal overlap is low throughout.}]{D:/face/artifacts/abliteration/gemma-3-1b-fp32_vs_gemma-3-1b-fp32-abl.png}
\caption{\textbf{Figure 3.} Abliteration on \protect\texttt{gemma-3-1b}
(fp32): every stock↔abliterated cosine (red) sits inside the
within-model self-cosine band (grey); emotion↔refusal overlap is low
throughout.}
\end{figure}

\textbf{7.3 Why: emotions are nearly orthogonal to the refusal
direction.} The null result has a direct geometric explanation. Each
emotion direction's overlap with the refusal direction is small in the
stock model --- \(|\cos(u_e, r)| \in [0.04, 0.10]\) --- and abliteration
leaves it small (Figure 3, lower). Because the emotion directions barely
project onto \(r\) to begin with, orthogonalizing the weights against
\(r\) removes almost nothing from them. Removing refusal, on this model,
is close to an \emph{orthogonal} edit with respect to the emotion
subspace.

\textbf{7.4 Summary and a caveat.} On \texttt{gemma-3-1b}, refusal
abliteration does \textbf{not} rotate the seven emotion directions
beyond finite-sample noise, and the reason is their low geometric
overlap with the refusal direction. The one secondary signal worth
flagging is a modest drop in \emph{separability} for a few axes even
though their direction is preserved (surprise AUC 0.99→0.92, frustration
1.00→0.88, confidence 0.72→0.61): abliteration may slightly blur how
cleanly some emotions are linearly decodable without reorienting them
--- a small effect at the edge of our sampling noise that larger prompt
sets should resolve. We read this 1B result as \emph{reassuring but
local}: on a small model with weakly entangled features, ``just removing
refusals'' is approximately an orthogonal edit. Whether that holds at
larger scale --- where emotion and refusal may share more geometry ---
is exactly what the scale sweep in §7.5 tests, and the answer turns out
to be \textbf{no}.

\begin{center}\rule{0.5\linewidth}{0.5pt}\end{center}

\textbf{7.5 Across scale, and a decisive control: the effect is
refusal-specific.} We repeated the experiment on the \texttt{gemma-3}
ladder in bf16 (Table 3, ``refusal'' columns). At 1B and 27B the emotion
directions are \textbf{preserved} (mean cos 0.87 / 0.96, ≤1/7 axes
moved, negligible AUC drop). At the \textbf{intermediate scales (4B,
12B)} we instead see a \textbf{large rotation} (mean cos 0.11 / 0.47,
7/7 axes ``moved'') accompanied by only a \emph{modest} separability
loss (mean AUC 0.93→0.82 at 4B, 0.97→0.91 at 12B). The emotions remain
well-separated after the edit --- this is a \textbf{re-orientation} of
the direction, not destruction of the concept.

Our first instinct was that this mid-scale rotation was a generic
side-effect of orthogonalizing weights (i.e.~the edit destabilizing the
model), especially since the emotion↔refusal \emph{linear} overlap is
small (\textbar cos\textbar{} ≈ 0.03). A \textbf{matched
random-direction ablation control} --- the \emph{identical}
weight-orthogonalization procedure applied to a random unit direction of
the same dimensionality --- refutes that (Table 3, ``random'' columns).
Random ablation leaves the emotion directions \textbf{essentially
untouched at every scale}: cos(stock, random-abl) = 0.95 / 0.99 / 1.00
for 1B/4B/12B, with 0/7 axes moved and \emph{no} AUC drop (and, as
expected, it does \textbf{not} reduce the refusal rate). So the
procedure itself is inert; it is \textbf{removing the refusal direction
specifically} that re-orients the emotion directions at 4B/12B. The
small \emph{linear} emotion↔refusal overlap does not bound this, because
abliteration is a weight edit whose effect propagates non-linearly
through the network --- it changes how emotions are \emph{computed}, not
merely subtracts a component.

The scientific reading: at intermediate scale, refusal and emotion
computation are \textbf{entangled} enough that excising refusal
reorganizes the emotion representation, whereas at the smallest scale
(weakly-represented refusal) and the largest (redundantly-represented
refusal, and a model that barely refused our probes) the effect is
small. This is a sharper and more surprising claim than §7.2's local
result, and it is \emph{only} legible because the control separates
``the concept we removed'' from ``the surgery we performed.'' It also
refines the safety message: ``just removing refusals'' is \textbf{not}
an orthogonal edit at every scale --- at mid-scale it measurably
restructures the model's affective representations, which our instrument
detects.

\begin{figure}
\centering
\includegraphics[width=0.95\linewidth,height=\textheight,keepaspectratio,alt={Figure 5. Refusal abliteration on gemma-3-4b (bf16): the emotion directions re-orient (red points fall well below the noise band) while separability is largely preserved --- an effect the random-direction control (Table 3) shows is refusal-specific.}]{D:/face/artifacts/abliteration/gemma-3-4b_vs_gemma-3-4b-abl.png}
\caption{\textbf{Figure 5.} Refusal abliteration on
\protect\texttt{gemma-3-4b} (bf16): the emotion directions re-orient
(red points fall well below the noise band) while separability is
largely preserved --- an effect the random-direction control (Table 3)
shows is refusal-specific.}
\end{figure}

\textbf{Table 3.} Refusal ablation vs.~a matched random-direction
control (\texttt{gemma-3}, bf16). Low cos(stock,abl) = large
re-orientation; the control isolates the \emph{refusal-specific} effect.

{\def\LTcaptype{none} % do not increment counter
\begin{longtable}[]{@{}
  >{\raggedright\arraybackslash}p{(\linewidth - 14\tabcolsep) * \real{0.1000}}
  >{\raggedright\arraybackslash}p{(\linewidth - 14\tabcolsep) * \real{0.1000}}
  >{\raggedleft\arraybackslash}p{(\linewidth - 14\tabcolsep) * \real{0.1333}}
  >{\centering\arraybackslash}p{(\linewidth - 14\tabcolsep) * \real{0.1333}}
  >{\raggedleft\arraybackslash}p{(\linewidth - 14\tabcolsep) * \real{0.1333}}
  >{\raggedleft\arraybackslash}p{(\linewidth - 14\tabcolsep) * \real{0.1333}}
  >{\centering\arraybackslash}p{(\linewidth - 14\tabcolsep) * \real{0.1333}}
  >{\raggedleft\arraybackslash}p{(\linewidth - 14\tabcolsep) * \real{0.1333}}@{}}
\toprule\noalign{}
\begin{minipage}[b]{\linewidth}\raggedright
model
\end{minipage} & \begin{minipage}[b]{\linewidth}\raggedright
refusal rate
\end{minipage} & \begin{minipage}[b]{\linewidth}\raggedleft
refusal: cos
\end{minipage} & \begin{minipage}[b]{\linewidth}\centering
mv/7
\end{minipage} & \begin{minipage}[b]{\linewidth}\raggedleft
ΔAUC
\end{minipage} & \begin{minipage}[b]{\linewidth}\raggedleft
random: cos
\end{minipage} & \begin{minipage}[b]{\linewidth}\centering
mv/7
\end{minipage} & \begin{minipage}[b]{\linewidth}\raggedleft
ΔAUC
\end{minipage} \\
\midrule\noalign{}
\endhead
\bottomrule\noalign{}
\endlastfoot
gemma-3-1b & 0.81 → 0.00 & 0.87 & 1/7 & +0.05 & 0.95 & 0/7 & +0.05 \\
gemma-3-4b & 0.44 → 0.00 & \textbf{0.11} & 7/7 & +0.11 & \textbf{0.99} &
0/7 & +0.01 \\
gemma-3-12b & 0.63 → 0.00 & \textbf{0.47} & 7/7 & +0.06 & \textbf{1.00}
& 0/7 & +0.01 \\
gemma-3-27b & 0.06 → 0.00 & 0.96 & 0/7 & +0.01 & --- & --- & --- \\
\end{longtable}
}

(Caveat retained: refusal \emph{rates} come from small stub sets + a
keyword detector and are unreliable --- 27B's 0.06 is implausibly low
--- so the \emph{quality} of the refusal direction varies across the
ladder; a standard benchmark, §10, would sharpen the mid-scale story
further. The random control, however, holds regardless of refusal-rate
calibration.)

\subsection{8. Discussion}\label{discussion}

That refusal abliteration re-orients emotion directions at some scales
but not others is informative in two ways. Practically, it bears on the
safety of ``uncensored'' open models: a weight edit marketed as
\emph{only} removing refusals demonstrably reshapes the model's
affective computation at mid-scale, and our instrument makes that
measurable rather than anecdotal --- with the random-direction control
(§7.5) showing the effect is genuinely tied to refusal, not to the
surgery. Scientifically, the scales where emotions move most (4B, 12B)
are where refusal and affect are most entangled --- a lead for
mechanistic follow-up, and a caution that the smallest and largest
models are \emph{not} representative of the middle.

We stress the limits of the claim. Cosine movement of a \emph{reading}
direction is evidence of geometric change, not yet of behavioural
change; establishing the latter (does an abliterated model's fear or
warmth behave differently, not just read differently?) is the next step.
This paper deliberately establishes \textbf{that} there is (or is not) a
measurable effect before asking \textbf{what} it does.

\begin{center}\rule{0.5\linewidth}{0.5pt}\end{center}

\subsection{9. Limitations}\label{limitations}

\begin{enumerate}
\def\labelenumi{(\roman{enumi})}
\tightlist
\item
  A 1B model has weaker, occasionally entangled emotion representations;
  steering has a narrow coherent range. (ii) Our shipped
  harmful/harmless sets are small stubs; the refusal direction and its
  overlaps should be re-estimated on a standard benchmark (AdvBench +
  Alpaca) at scale. (iii) Difference-of-means treats each emotion as one
  direction; genuinely multi-dimensional concepts are under-modelled.
  (iv) The abliteration-perturbation result is measured on reading
  directions (cosine of the readout direction), not yet on behaviour;
  the matched random-direction control (§7.5) rules out a weight-surgery
  artifact, but a behavioural probe of the re-oriented mid-scale models
  is future work. (v) The bipolar confidence axis and the calibration
  baseline are model- and prompt-set-specific.
\end{enumerate}

\begin{center}\rule{0.5\linewidth}{0.5pt}\end{center}

\subsection{10. Future work}\label{future-work}

\textbf{Scale.} Repeat the full experiment on Gemma-4-E4B (stock
vs.~abliterated twin) on a 24 GB GPU; the pipeline is model-agnostic and
the driver is parameterized. \textbf{Controls.} The matched
random-direction ablation (§7.5) establishes specificity; next is
dose-response (varying strength / rank) and a \emph{behavioural} readout
of the re-oriented mid-scale models --- does reorganized affect change
what they do? \textbf{Discovery.} Train sparse autoencoders on both
twins and ask whether the \emph{features} --- not just the named axes
--- shift, and whether new affective features appear.
\textbf{Mechanism.} For the emotions that move, trace \emph{what
changes}: which behaviours, on which inputs, and whether the change is
mediated by the removed refusal direction. This is the journey the
present paper begins.

\begin{center}\rule{0.5\linewidth}{0.5pt}\end{center}

\subsection{11. Reproducibility
statement}\label{reproducibility-statement}

The toolkit, prompt sets, fitting/calibration code, steering and
abliteration implementations, the experiment driver, the figures, and
the per-run JSON results bundles (with software versions) are released
under the MIT license at
\href{https://github.com/newsbubbles/faced}{github.com/newsbubbles/faced}.
The abliteration is a closed-form weight edit with a seeded,
deterministic direction; the statistical test is a seeded bootstrap. A
single command reproduces each figure from the released prompt sets (see
the repository \texttt{README.md} and \texttt{WALKTHROUGH.md}). The
large binary artifacts --- the trained SAE and the fitted direction
tensors --- are regenerable from that code and are hosted separately
rather than committed, so the repository stays light.

\begin{center}\rule{0.5\linewidth}{0.5pt}\end{center}

\subsection{Figure index}\label{figure-index}

Generated artifacts (regenerate with the commands in \texttt{README.md}
/ \texttt{WALKTHROUGH.md}):

\begin{itemize}
\tightlist
\item
  \textbf{Figure 1} --- emotion meters over generation time:
  \texttt{artifacts/emotions\_over\_time.png}
\item
  \textbf{Figure 2} --- steering:
  \texttt{artifacts/steering\_eval.gemma-3-1b.txt} (α-sweep,
  suppression, confusion matrix) and the peak-emotion face gallery
  \texttt{artifacts/faces/gallery.html}
\item
  \textbf{Figure 3} --- abliteration effect (cosine vs.~noise floor;
  refusal overlap), local fp32 1B:
  \texttt{artifacts/abliteration/gemma-3-1b-fp32\_vs\_gemma-3-1b-fp32-abl.png}
\item
  \textbf{Figure 4} --- signal cleanliness vs.~scale (AUC / d′ /
  self-cosine heatmaps): \texttt{artifacts/scaling/cleanliness.png};
  per-model abliteration in
  \texttt{artifacts/abliteration/gemma-3-*\_vs\_*-abl.png}.
\item
  Raw results + provenance: \texttt{artifacts/scaling/cleanliness.json},
  \texttt{artifacts/abliteration/*\_vs\_*-abl.json} (bf16 ladder) + the
  fp32 1B json, plus per-model abliteration manifests.
\end{itemize}

\bibliography{D:\textbackslash face\textbackslash paper\textbackslash references.bib}

\end{document}
